\documentclass{article}
\usepackage{graphicx} % Required for inserting images
\usepackage{amsmath}
  \RequirePackage[french]{babel}
  \RequirePackage[T1]{fontenc}
  \RequirePackage[utf8]{inputenc}

\title{STT-2902 - Modélisation statistique}
\author{Julien Miron}


\begin{document}

\section*{Équipe 5 – Régression linéaire}

\textbf{Session :} Automne 2026

\textbf{Style imposé : séance de révision d'examen.} Votre présentation doit prendre la forme d'une séance de révision en vue de l'examen : vous préparez des questions de type examen sur la régression linéaire, les faites résoudre à l'auditoire, puis présentez des solutions détaillées. Ce style est obligatoire.

\vspace{0.5cm}
\textbf{Objectif :} Consolider la compréhension de la régression linéaire simple : estimation de la droite, interprétation des coefficients, qualité de l'ajustement et utilisation du modèle pour la prédiction.

\vspace{0.5cm}
\textbf{Déroulement imposé} :
\begin{itemize}
    \item Un bref rappel des notions essentielles (au plus 10 minutes) : modèle, pente, ordonnée à l'origine, $R^2$, prédiction;
    \item Une série de questions de type examen que l'auditoire résout individuellement ou en petit groupe, avec un temps limité par question;
    \item Une correction détaillée de chaque question, en insistant sur les erreurs fréquentes;
    \item Une liste de \og pièges classiques \fg{} à éviter à l'examen.
\end{itemize}

\vspace{0.5cm}
\textbf{Requis} :
\begin{enumerate}
    \item Préparer au moins 4 questions de type examen, dont au moins une question d'interprétation des coefficients et une question de prédiction;
    \item Inclure une question portant sur l'effet de valeurs aberrantes sur la droite de régression;
    \item Inclure une question où la relation entre les deux variables n'est pas linéaire, et expliquer pourquoi le modèle linéaire ne convient pas;
    \item Aborder les dangers de l'extrapolation dans au moins une question ou sa correction;
    \item Fournir un corrigé écrit complet remis à l'enseignant avec le matériel de présentation.
\end{enumerate}

\end{document}
